% Source: https://artofproblemsolving.com/community/c3h1463899_functions_question
% Stable identity: aops-topic-1463899-q01
The symbol $\sqrt{x}$ means the \emph{principal} (non-negative) square root of $x$: the unique number $y\ge0$ with $y^2=x$. For instance $\sqrt{4}=2$, not $-2$. So each input $x\ge0$ has exactly one output, and $f$ is a function. The other square root of $x$ is written $-\sqrt{x}$.

The equation $y^2=x$, on the other hand, does \textbf{not} define $y$ as a function of $x$: for $x=4$ both $y=2$ and $y=-2$ satisfy it. Its graph is a sideways parabola, which fails the vertical line test; its upper half is the graph of $y=\sqrt{x}$ and its lower half is the graph of $y=-\sqrt{x}$, and each half is a function.
